% Draft for papers/parity-barrier/main.tex, thread 2. Not yet inserted.
% Source notes: notes/sources-mr16.md (arXiv:1501.04585v4, SHA-256 ec546fdf...9277).
\subsection{Multiplicative functions in short intervals}\label{subsec:mr16}
\lbl{Literature} \cite{MR16} (arXiv:1501.04585v4; PDF page $=$ printed page).
Theorem~1 (pp.~1--2): let $f:\mathbb N\to[-1,1]$ be multiplicative. There are
absolute constants $C,C'>1$ such that for any $2\le h\le X$ and $\delta>0$,
\[
\Bigl|\frac1h\sum_{x\le n\le x+h}f(n)-\frac1X\sum_{X\le n\le2X}f(n)\Bigr|
\le\delta+C'\frac{\log\log h}{\log h}
\]
for all but at most
$CX\bigl((\log h)^{1/3}\delta^{-2}h^{-\delta/25}+\delta^{-2}(\log X)^{-1/50}\bigr)$
integers $x\in[X,2X]$; one can take $C'=20000$. So $f$ is real-valued,
$1$-bounded and multiplicative (not necessarily completely multiplicative),
and the statement is non-trivial as soon as $h\to\infty$. Consequences (pp.~3--5):
$\frac1X\bigl|\sum_{n\le X}\lambda(n)\lambda(n+h)\bigr|\le1-\delta(h)$ for
every $h\ge1$ (Corollary~2), and for $f:\mathbb N\to\mathbb R$ multiplicative,
negative somewhere and non-zero on a positive proportion of integers, almost
every $[x,x+\psi(x)]$ with $\psi(x)\to\infty$ contains a sign change of $f$
(Corollary~4). In particular, every $[n,n+C\sqrt n]$ contains integers
with an even number and with an odd number of prime factors (p.~5, from
Corollary~5). The structure of the proof is set out on pp.~6--11:
\begin{enumerate}
\item \emph{Typical integers.} Restrict to the set $S$ of $n\in[X,2X]$
  with a prime factor in each of a chain of ranges $[P_j,Q_j]$, where
  $Q_1\le\exp(\sqrt{\log X})$. Most $n$ lie in $S$, and the rest are removed by
  the fundamental lemma of the sieve (pp.~6, 30--31).
\item \emph{Parseval.} The variance over $x$ reduces to the mean square
  of $\sum_{n\in S}f(n)n^{-1-it}$ over $(\log X)^{1/15}\le|t|\le X/h$
  (p.~7, (5); Lemma~14, p.~21).
\item \emph{Factorisation.} Where the prime polynomial
  $\sum_{P_j\le p\le Q_j}f(p)p^{-1-it}$ is small, ``a variant of Ramar\'e's
  identity'' factors out this polynomial. A pointwise bound for it is combined
  with the mean value theorem for the cofactor, and moments are used for $j>1$
  (pp.~9--10; Lemma~12, p.~19).
\item \emph{Large values and Hal\'asz.} The remaining sparse set of $t$ is
  handled with Hal\'asz-type large-value inequalities (Lemmas~9 and~11,
  pp.~16--17) and Hal\'asz's theorem. That theorem makes the long average the
  benchmark and uses that $f$ is real (pp.~10--13).
\end{enumerate}
The authors say that the result fails for complex $f$ (for example
$f(p)=p^{it}$), and that it cannot hold in all intervals of length
$\asymp\sqrt X$, because ``the large primes factors'' are ``completely arbitrary
for general $f$'' (pp.~2, 5). They note that Chowla's conjecture ``is believed
to be at least as deep as the twin prime conjecture'' (p.~3). The word
``parity'' does not occur in the paper.

\begin{remark}[Why this does not transfer to twin primes]\label{rem:mr16-transfer}
\lbl{Heuristic} (our reading, not a statement of \cite{MR16}). Every step
above applies to a single function $f(n)$ that is multiplicative, real and
bounded by~$1$. Step~(3) needs $f(pm)=f(p)f(m)$ (p.~9). The function
$n\mapsto\lambda(n)\lambda(n+2)$ is not multiplicative, and
$\Lambda(n+2)$ is neither bounded nor multiplicative in~$n$. So
Theorem~1 says nothing about either of them directly. For
$\lambda(n)\lambda(n+h)$ the paper gets only $1-\delta(h)$, through sign
changes (pp.~37--38). The two-point results we use come from a different
argument, \cite[Theorem~1.2]{Tao16}. It holds only with logarithmic averaging,
and its proof uses $\lambda(pn)=-\lambda(n)$ at small primes in a new
bilinear estimate (Section~\ref{sec:parity}); we have not checked here how it
uses \cite{MR16}. Twin primes need more than a short-interval average of
one function: they need a correlation $\sum_n\Lambda(n)\Lambda(n+2)$ of
$\Lambda$ with a shift of itself. Nothing in \cite{MR16} addresses that
correlation (\lbl{Open}).
\end{remark}

% bibitem for the thebibliography environment:
\bibitem{MR16} K.~Matom\"aki and M.~Radziwi\l\l, \emph{Multiplicative functions in short intervals}, Ann. of Math. (2) \textbf{183} (2016) (journal version not consulted; reference as supplied, not printed in the PDF); arXiv:1501.04585v4 (15 Oct 2017; read: pp.~1--19, 21--24, 29--31, 36--38).
